% Responsable(s): por asignar
% Objetivo: sintetizar tendencias, enfoques actuales y hallazgos recurrentes
% identificados en los artículos seleccionados.

% Redacción pendiente.
En conjunto, estos cuatro artículos cubren tanto memoria principal ya consolidada ---DDR4,
DDR5, HBM, HBM2, HBM3--- como tecnologías más nuevas de expansión y caché de memoria
---CXL, eDRAM, MRAM---, todos con datos medidos en hardware o en simulación detallada, y
todos publicados entre 2023 y 2025. Esa tensión entre ancho de banda y latencia que aparece una
y otra vez, junto con el costo que implica salir de los canales DIMM convencionales, es la
evidencia que este bloque aporta al análisis comparativo del grupo.
De manera complementaria, Eyerman et al. proponen analizar el comportamiento
de la memoria DRAM mediante \textit{bandwidth stacks} y \textit{latency stacks},
con el objetivo de identificar las causas que limitan el aprovechamiento del
ancho de banda y aumentan la latencia. Su trabajo muestra que el rendimiento
real de DRAM está condicionado por restricciones temporales, paralelismo interno
y esperas en el controlador de memoria, por lo que las especificaciones teóricas
no representan necesariamente el comportamiento observado en una aplicación
\cite{eyerman2022dram}.